\section{Experimental Setup}\label{sec:setup}

All evaluation ran on one laptop with Docker, giving the cluster about 14~GB of memory. The cluster is a local Kubernetes cluster created with \texttt{kind} \citep{tool_kind}, which runs each Kubernetes node as a Docker container. The experiment therefore exercises the three enabling technologies of cloud-native systems directly: containers (every node and every workload), orchestration (Kubernetes schedules, restarts and replaces the workloads, which D5 and D6 measure) and isolation (every configuration gets its own namespace). No public cloud account is needed, so the setup can be reproduced at no cost. Table~\ref{tab:env} lists the pinned versions; the scanners use rule sets and schemas downloaded once and reused offline, so every configuration is judged by the same rules.

\begin{table}[t]
\centering\small
\caption{Pinned evaluation environment (\texttt{tools/versions.lock}).}\label{tab:env}
\begin{tabular}{ll}
\toprule
Component & Version \\ \midrule
Docker Engine & 29.1.3 \\
kind & v0.33.0 \\
Kubernetes (node image), kubectl & v1.37.0 \\
Cluster & 1 control plane + 3 workers \\
kubeconform & v0.8.0 \\
Kubescape & v4.0.16 \\
Trivy & v0.76.0 \\
kube-score & v1.20.0 \\
Polaris & v10.2.5 \\
KubeLinter & v0.8.3 \\
\bottomrule
\end{tabular}
\end{table}

\subsection{Evaluation pipeline and environment}\label{sec:method-pipeline}

Each configuration passes through the pipeline in a fixed order (Algorithm~\ref{alg:evaluate}): parsing, schema validation, the server-side dry-run, the D2/D3 scanners, the D4 checks, deployment and smoke test (D5), and fault injection (D6). Failures are recorded as data. A parse failure marks every later stage as skipped. A schema failure does not block the static scanners or D4, which read the manifest text. D6 runs when the configuration is fully Ready. Otherwise, it runs when its smoke test passes, because an application can serve correctly while an auxiliary workload (for example, a load generator) is not Ready. We also report D6 for fully Ready configurations only.

Every configuration is deployed into its own namespace, and we recreate the cluster between projects. Every image comes from a local registry mirror filled before deployment, so network latency does not confound time to Ready. Two deployment-time adjustments supply facts that the bundles do not contain. Both are applied identically to every source and recorded per configuration, and no human manifest ever needs them. First, in some projects the exact image path appears only in unbundled build scripts. There, a reference that cannot be pulled as written but whose repository name identifies a documented component is served that component's canonical image. A reference that identifies no documented component fails. Second, an image may enforce a start-up prerequisite that the bundle never states. We supply it when the configuration does not provide it. These prerequisites are two environment variables in Online Boutique (one its front end requires, one that disables a cloud profiler in two services), and an empty Secret with the fixed name that the Kubernetes Dashboard expects. Permissions are never added. Section~\ref{sec:threats} bounds the effect of both adjustments.

The cluster has one control-plane node and three worker nodes. So configurations that use pod anti-affinity or a topology spread across three replicas can schedule. One setting differs from \texttt{kind}'s defaults. The container runtime limits each container to 1,048,576 open files (Docker's default). We set this because some older images size internal tables from that limit and could otherwise exhaust the host's memory. Image tags are resolved when the mirror is filled, at evaluation time. A \texttt{LoadBalancer} Service is treated as valid but pending. An Ingress is created but not exercised. PersistentVolumeClaims use \texttt{kind}'s local-path StorageClass. So a configuration that requests a StorageClass absent from this environment fails, and the reason is recorded.

\subsection{Implementation: data structures, algorithms and design patterns}\label{sec:setup-impl}

The pipeline is written in Python. \texttt{eval\_runner.py} runs D1, D4 and the order of stages; \texttt{d2d3.py} runs the static scanners and the checklist; \texttt{d5d6.py} deploys, tests and injects faults. The external tools are called as command-line programs.

\paragraph{Data structures.} Three structures carry the whole study.
\begin{itemize}
\item The \emph{result record}: one JSON document per configuration, keyed by project, source, prompt and sample. It has one field per dimension (\texttt{d1\_*}, \texttt{d2}, \ldots, \texttt{d6}). Every stage writes its measurements and, when it does not run, a reason such as \texttt{SKIPPED(parse\_error)}, so no configuration is ever dropped. All tables and figures are computed from these 205 records.
\item The \emph{project specifications}: a declarative smoke test per project (\texttt{smoke.json}: a list of HTTP requests or data-store commands, each with its expected status, body or pattern) and a requirement script (\texttt{requirements.py}) written as calls to a shared \texttt{Checker} library, such as ``image deployed'', ``Service named \ldots'' or ``volume persistent''. Adding a project means adding data, not changing the pipeline.
\item The \emph{evidence tables}: the root-cause labels (\texttt{d5\_root\_causes.csv}: configuration, category, mechanism, evidence level, source) and the pinned environment (\texttt{versions.lock}).
\end{itemize}

\paragraph{Algorithms.} Algorithm~\ref{alg:evaluate} shows how one configuration is evaluated. Three steps need care. The YAML loader rejects duplicate mapping keys, which the default loader silently merges. Readiness is detected with a stability window: the cluster is polled every 3~s, and a configuration counts as Ready only when every workload has held its desired Ready replicas for 20~s, within a 300~s deadline, so that a pod that crashes right after becoming Ready is not counted. Downtime in D6 is the number of failed probes divided by the probing rate.

\begin{algorithm}[t]
\caption{Evaluating one configuration $c$ of project $p$.}\label{alg:evaluate}
\small
\begin{algorithmic}[1]
\State $r \gets$ new result record for $c$
\State $(docs, text, err) \gets \Call{StrictParse}{c}$ \Comment{duplicate keys rejected}
\If{$err \neq \varnothing$}
  \State mark D1--D6 of $r$ as skipped; \Return $r$
\EndIf
\State $r.\mathrm{D1} \gets \Call{Kubeconform}{text}$, \Call{ServerDryRun}{$docs$} \Comment{object by object}
\State $r.\mathrm{D2}, r.\mathrm{D3} \gets \Call{Scanners}{docs, text}$, \Call{PodSecurityDryRun}{$docs$}
\State $r.\mathrm{D4} \gets \Call{Requirements}{p, docs}$
\State $ns \gets$ fresh namespace; serve images from the local mirror; apply $docs$ to $ns$
\State $t_0 \gets now$; $since \gets \bot$
\While{$now - t_0 < 300$\,s}
  \If{every workload has its desired Ready replicas}
    \State $since \gets since$ \textbf{or} $now$
    \If{$now - since \geq 20$\,s}
      \State \textbf{break}
    \EndIf
  \Else
    \State $since \gets \bot$
  \EndIf
  \State wait 3\,s
\EndWhile
\State $r.\mathrm{ready} \gets (since \neq \bot)$; \textbf{if not} $r.\mathrm{ready}$ \textbf{then} capture diagnostics
\State $r.\mathrm{smoke} \gets \Call{SmokeTest}{p, ns}$ \Comment{client placed where NetworkPolicies admit it}
\If{$r.\mathrm{ready}$ \textbf{or} $r.\mathrm{smoke}$ passed}
  \State $r.\mathrm{D6} \gets$ 3 pod deletes and 3 container crashes under a 5\,Hz prober
\EndIf
\State delete $ns$ \Comment{always, even after an error}
\State \Return $r$
\end{algorithmic}
\end{algorithm}

\paragraph{Design patterns.} The pipeline follows a small set of familiar patterns. It is a \emph{pipeline} of stages with a fixed order, in which a failed stage is recorded and later stages are skipped or still run according to explicit rules. Each external tool sits behind an \emph{adapter}: one function per tool (\texttt{kubescape()}, \texttt{trivy()}, \texttt{polaris()} and so on) that runs it, keeps its raw output and converts it into the common record format, so tools can be added or replaced without touching the rest. Project-specific behaviour is \emph{data-driven}, through the declarative smoke tests and requirement checks above. Cluster resources are released in a \emph{try/finally} block, so a failed measurement never leaves a namespace behind. And the environment is \emph{pinned}: tool versions, scanner rules, schemas and images are fixed once and reused, so a re-run gives the same static verdicts.

\subsection{Statistical analysis}\label{sec:method-analysis}

We report medians for scores and percentages for pass/fail outcomes. Because the three samples of a (model, project, prompt) cell share their input, and the projects differ widely, our tests treat the project or the (model, project) cell as the unit wherever possible. We use rank-based tests: Kruskal--Wallis and Fisher's exact test to compare sources (RQ1), a Friedman test blocked on project to compare models (RQ3), a logistic regression to relate static scores to deployment success (RQ4), and a Wilcoxon signed-rank test to compare the two prompts within each (model, project) cell (RQ5). $P$-values are Holm-adjusted for multiple comparisons. Appendix~\ref{app:stats} describes every test, the effect sizes and the robustness checks, and reports their results.
